\chapter{Evaluation Artifacts and Reproduction}
\label{app:commands}

\section{Official Experiment Commands}

The following commands are executed from the backend directory in Windows CMD.
The same work directory must be used for all stages.

\begin{verbatim}
cd /d E:\Thesis_all\be
set PYTHON_RUNNER=.venv\Scripts\python.exe

%PYTHON_RUNNER% evaluation\run_openai_experiment.py ^
  --stage generate --users 5 --conversations 2 --turns 10 ^
  --work-dir evaluation\runs\v25_official_100

%PYTHON_RUNNER% evaluation\run_openai_experiment.py ^
  --stage pipeline --limit 100 ^
  --work-dir evaluation\runs\v25_official_100

%PYTHON_RUNNER% evaluation\run_openai_experiment.py ^
  --stage judge --limit 100 ^
  --work-dir evaluation\runs\v25_official_100

%PYTHON_RUNNER% evaluation\run_openai_experiment.py ^
  --stage export --limit 100 ^
  --work-dir evaluation\runs\v25_official_100
\end{verbatim}

The pipeline and judge stages are resumable. If the judge is started while the
pipeline is still running, it evaluates only the cases present in the pipeline
trace when the judge process starts. Running the judge again after all pipeline
cases are complete reuses compatible judged cases and processes the remainder.

\section{Artifact Files}

\begin{table}[H]
\centering
\caption{Official experiment artifacts}
\begin{tabular}{p{0.40\textwidth} p{0.50\textwidth}}
\hline
\textbf{File} & \textbf{Purpose} \\
\hline
\texttt{01\_generated\_dataset.json} & Profiles, memories, conversations,
queries, and draft references \\
\texttt{02\_pipeline\_traces.json} & Intermediate stages, evidence, answer,
timings, and pipeline telemetry \\
\texttt{03\_judged\_traces.json} & Retrieval grades, final-answer scores, and
judge call metadata \\
\texttt{04\_thesis\_dataset.json} & Validated dataset used for deterministic
metric calculation \\
\texttt{05\_metric\_summary.json} & Aggregate metrics, usage, latency, cost,
models, prompts, and corpus metadata \\
\hline
\end{tabular}
\end{table}

\section{Query Taxonomy}

The supported query intents are:

\begin{itemize}
    \item itinerary;
    \item recommendation;
    \item attraction search;
    \item accommodation search;
    \item food search;
    \item transport;
    \item event search; and
    \item travel information.
\end{itemize}

The supported operations are lookup, recommend, compare, plan, and explain.
Intent and operation should be annotated independently.

\section{Retrieval Relevance Rubric}

\begin{table}[H]
\centering
\caption{Chunk relevance grading rubric}
\begin{tabular}{cl}
\hline
\textbf{Grade} & \textbf{Interpretation} \\
\hline
0 & Unrelated to the query \\
1 & General background but not sufficient for an answer \\
2 & Useful for one part or aspect of the query \\
3 & Directly useful evidence for the query \\
\hline
\end{tabular}
\end{table}

Grades 2 and 3 are considered relevant for Precision@5. All four grades are
used by nDCG@5.

\section{Final-Answer Judge Rubric}

Each dimension is scored from 1 to 5. A score of 1 represents a major failure,
3 represents a mixed or partially adequate response, and 5 represents strong
performance for that dimension. Correctness and faithfulness must be judged
against the supplied evidence rather than the evaluator's unsupported external
knowledge. Completeness must account for answer readiness: an honest partial
answer may be faithful even when it cannot fully satisfy the query.

\section{Human Review Checklist}

For every selected review case, the reviewer should check:

\begin{enumerate}
    \item whether the query is coherent in its conversation;
    \item whether intent and operation labels are defensible;
    \item whether explicit constraints exclude inferred profile preferences;
    \item whether retrieval facets follow the controlled vocabulary;
    \item whether broad regions are not mislabeled as provinces or cities;
    \item whether negative constraints remain negative;
    \item whether relevance grades follow the 0--3 rubric;
    \item whether factual answer claims are supported by selected evidence;
    \item whether relevant preferences are applied without forcing irrelevant
    personalization;
    \item whether missing evidence is acknowledged; and
    \item whether the judge rationale is consistent with its numerical scores.
\end{enumerate}

\section{Results Completion Checklist}

Before submitting the thesis:

\begin{itemize}
    \item replace all \textit{Pending official run} placeholders;
    \item verify that every reported value matches the exported summary or a
    documented analysis script;
    \item add the official corpus counts and fingerprint;
    \item record the number and method of human-reviewed cases;
    \item insert figures generated from the official dataset only;
    \item update the conclusion with evidence-based aggregate results;
    \item confirm that API keys and user-sensitive information are absent; and
    \item archive the prompts, schemas, code commit, and run artifacts.
\end{itemize}
